% исправления 27.05.2005
\section{Переименование переменных}
        \label{renaming-variables}
В этом разделе мы попытаемся аккуратно разобраться с простым
вопросом о том, почему и как можно переименовывать связанные
переменные, не меняя смысла формул. Мы уже видели, что формулы
$\forall x\,A(x)$ и $\forall y\,A(y)$ доказуемо эквивалентны
(с.~\pageref{changing-variables-example}), то есть их
эквивалентность доказуема в исчислении предикатов. Сейчас мы
хотим доказать общее утверждение об этом.

Корректная формулировка утверждения о переименовании переменных
требует осторожности. Например, нельзя сказать, что формула
$\forall\xi\, \varphi$ всегда эквивалентна
$\forall\eta\,\varphi(\eta/\xi)$. Прежде всего, подстановка может
быть некорректной, как в случае формул
        $$
\forall \xi \forall\eta\,  B(\xi,\eta) \quad\text{и}\quad
\forall \eta \forall\eta\,  B(\eta,\eta)
        $$
(легко понять, что эти формулы не эквивалентны).
Но даже если подстановка корректна, формулы могут не быть
эквивалентными, как в случае формул
        $$
\forall \xi\, B(\xi,\eta) \quad\text{и}\quad
\forall \eta\, B(\eta,\eta).
        $$
Как же сформулировать утверждение правильно?

Нагляднее всего, видимо, сделать так. Давайте заключим в рамку
все связанные вхождения всех переменных (в том числе вхождения
после кванторов). После этого соединим линиями переменную после
квантора и все её вхождения, связанные именно этим вхождением
квантора. Свободные вхождения переменных остаются при этом без
рамок. Получится что-то вроде
$$
\unitlength=1mm
        %
\begin{picture}(90,10)
        %
\put(8,6){$\forall \framevar{x}\,(\exists\framevar{z}\,
        B(\framevar{x},\framevar{z}) \land C(\framevar{x}))
        \land D(x,y,z).$}

\put(13,2){\line(0,1){2}}
\put(34.1,2){\line(0,1){2}}
\put(58,2){\line(0,1){2}}
\put(13,2){\line(1,0){45}}
\put(23,0){\line(0,1){4}}
\put(42,0){\line(0,1){4}}
\put(23,0){\line(1,0){19}}
            %
\end{picture}
$$
Если теперь стереть переменные внутри рамок, останется
\emph{схема}%
\index{Схема формулы|defin}%
\index{Формула!схема|defin}
формулы, которая содержит всю существенную
информацию о ней. Будем называть две формулы \emph{подобными}%
\index{Подобные формулы|defin}%
\index{Формулы!подобные|defin}
(отличающимися лишь именами связанных переменных), если они
имеют одну и ту же схему.

\begin{theorem}[о переименовании связанных переменных]
\index{Теорема!о переименовании связанных переменных|defin}%
        \label{renaming-variables-theorem}%
Подобные формулы доказуемо эквивалентны.
        %
\end{theorem}

\begin{proof}
        %
Докажем две простые леммы.

\textsf{Лемма~1}. Если формула $\varphi$ не содержит переменной~$\eta$
(ни связанно, ни свободно), то формулы
        $$
\forall\xi\,\varphi \quad \text{и}\quad \forall\eta\,\varphi(\eta/\xi)
        $$
доказуемо эквивалентны.

Доказательство. В самом деле, подстановка корректна, так как в~$\varphi$
нет кванторов по~$\eta$. Поэтому выводима формула
        $$
\forall\xi\,\varphi \to \varphi(\eta/\xi).
        $$
Левая часть её не содержит переменной $\eta$, поэтому по правилу
Бернайса можно вывести
        $$
\forall\xi\, \varphi \to \forall\eta\,\varphi(\eta/\xi).
        $$

В обратную сторону:
подстановка $\xi$ вместо~$\eta$ в формулу $\varphi(\eta/\xi)$
корректна (поскольку $\eta$ была подставлена вместо свободных
вхождений~$\xi$, при обратной подстановке переменная~$\xi$ не
попадёт в область действия кванторов по ней) и даёт формулу $\varphi$.
Поэтому формула
        $$
\forall\eta\,\varphi(\eta/\xi)\to \varphi(\eta/\xi)(\xi/\eta)
        $$
или, что то же самое,
        $$
\forall\eta\,\varphi(\eta/\xi)\to \varphi
        $$
является аксиомой. Осталось применить правило Бернайса, заметив,
что в левую часть переменная $\xi$ свободно не входит (все свободные
вхождения были заменены на~$\eta$). Лемма~1 доказана.

Аналогичное утверждение для квантора существования доказывается
точно так же.

\textsf{Лемма~2.} Замена подформулы на доказуемо эквивалентную
даёт доказуемо эквивалентную формулу.
     \label{changing-parts}%

Доказательство. Как мы видели на
с.~\pageref{adding-quantifiers}, доказуемая эквивалентность
сохраняется после навешивания квантора: если
$\alpha\hm\simeq\alpha'$, то
$\forall\xi\,\alpha\hm\simeq\forall\xi\,\alpha'$ и
$\exists\xi\,\alpha\hm\simeq\exists\xi\,\alpha'$ (символ
$\simeq$ здесь обозначает доказуемую эквивалентность).
Кроме того, из $\alpha\hm\simeq\alpha'$ и $\beta\hm\simeq\beta'$
следует, что ${(\alpha\land\beta)}\hm\simeq{(\alpha'\land\beta')}$,
             ${(\alpha\lor\beta)}\hm\simeq{(\alpha'\lor\beta')}$,
             ${(\alpha\to\beta)}\hm\simeq{(\alpha'\to\beta')}$ и
             $\lnot\alpha\hm\simeq\lnot\alpha'$.
(В~этом легко убедиться, написав подходящие пропозициональные
тавтологии.)

Теперь утверждение леммы легко доказать индукцией (начав с заменённой
подформулы и рассматривая всё более длинные части формулы).
Лемма~2 доказана.

Леммы~1 и~2 позволяют нам заменять переменные внутри рамок схемы
на новые (ранее не использованные) переменные, получая доказуемо
эквивалентную и подобную исходной формулу. Такими заменами можно
из двух подобных формул получить третью, используя для замены
одни и те же переменные. При этом обе исходные формулы доказуемо
эквивалентны третьей, а значит, и друг~другу.

(Использование третьей формулы существенно: нельзя
преобразовать первую формулу сразу во вторую, так как при замене
переменных в рамках может не выполняться условие леммы~1.)
        %
\end{proof}

Аккуратное обращение со связанными и свободными переменными\т
традиционная головная боль авторов учебников по логике.
Наиболее радикальный подход\т вообще изгнать связанные
переменные, заменив их квадратиками со связями между ними. Тогда
при подстановке можно ни о чём не заботиться. Зато формулы
перестают быть последовательностями символов, а становятся
объектами со сложной структурой. (Этот подход использован в
книге Бурбаки%
\glossary{\Бурбаки}
\emph{Теория множеств}~\cite{bourbaki}.)

Менее радикальный вариант состоит в том, чтобы разделить
переменные на два типа\т свободные и связанные. Так делается,
например, в классической книге Гильберта%
\glossary{\Гильберт}
и Бернайса%
\glossary{\Бернайс}
\emph{Основания математики}~\cite{hilbert-bernays}. Тогда можно
смело подставлять терм вместо свободной переменной, зато при
навешивании квантора надо заменять свободную переменную на
связанную.

Ещё один вариант\т договориться, что при подстановке терма
вместо свободной переменных автоматически происходит
переименование связанных переменных, создающих коллизии.

Всё это, конечно, мелочи\т но досадные, особенно если стремиться
к краткости, ясности и наглядности. Следы мучительных раздумий
на подобные темы видны в примечании на с.~101\ч 102 книги
Клини%
\glossary{\Клини}%
~\emph{Математическая
логика}~\cite{kleene-logic}: \лк
Гильберт и Бернайс~\hole и другие авторы
используют для обозначения свободных и связанных переменных
разные буквы~\hole Мы следовали этому правилу в
течение десятилетия~\hole Сейчас же мы твёрдо
убеждены, что использование единого списка переменных для
свободных и замкнутых вхождений даёт небольшое, но
чувствительное преимущество\пк.

С этим связан и другой выбор: как определять истинность формул.
Тут есть две возможности: можно определять истинность формулы на
оценке (при данных значениях параметров), а можно говорить
только о замкнутых формулах, вводя константы для всех элементов
интерпретации. И тот, и другой способы имеют недостатки, а выбор
нами первого подхода\т результат не единодушия авторов, а
волевого решения.

\section{Предварённая нормальная форма}
        \label{prenex-normal-form}%

Говорят, что формула находится в \emph{предварённой нормальной
форме},%
\index{Предварённая нормальная форма|defin}%
\index{Нормальная форма!предварённая|defin}
если все кванторы в ней вынесены налево, то есть 
она имеет вид
        $$
Q_1 \xi_1 \ldots Q_k \xi_k \varphi,
        $$
в котором $Q_1,\dots,Q_k$\т кванторы (всеобщности или существования),
$\xi_1,\dots,\xi_k$\т переменные, а $\varphi$\т бескванторная
формула. Эта формула может иметь параметры (если
формула~$\varphi$ имеет параметры, отличные от
$\xi_1,\dots,\xi_k$).

Основной результат этого раздела гласит, что всякая формула
(доказуемо) эквивалентна некоторой формуле в предварённой
нормальной форме (\emph{предварённой формуле}).%
\index{Предварённая формула|defin}%
\index{Формула!предварённая|defin}
Мы докажем его,
одновременно построив некоторую классификацию формул (в каком\д
то смысле отражающую их \лк логическую сложность\пк). В качестве
меры сложности можно было бы взять число кванторов в
предварённой нормальной форме. Но правильнее учитывать число
групп кванторов (считая одноимённые рядом стоящие кванторы за
один).

Говорят, что предварённая формула является $\Sigma_n$\д
формулой, если её кванторная приставка содержит $n$ групп
кванторов, причём первыми стоят кванторы существования. Если
первыми стоят кванторы всеобщности, говорят о классе~$\Pi_n$.
(Аналогичные обозначения используются в теории алгоритмов для
классификации арифметических множеств,
см.~\cite{computable-functions}).

Скажем, формула ${\forall x \exists y \exists z\forall u\,
(A(x,u,z)\to B(z,u))}$ принадлежит классу $\Pi_3$, формула
$\exists u\forall v\,C(u,v)$ принадлежит классу~$\Sigma_2$, а
формула ${\forall x\, (A(x) \to \exists y B(x,y))}$ вообще не
находится в предварённой нормальной форме.

\begin{problem}
        %
Укажите формулу в предварённой нормальной форме, доказуемо
эквивалентную последней из перечисленных формул.
        %
\end{problem}

Нас интересует, что происходит с измеряемой таким образом \лк
логической сложностью\пк\ формулы при логических операциях.
Начнём с совсем простых наблюдений.
        %
\begin{itemize}
        \item
Всякая формула из класса $\Sigma_n$ или $\Pi_n$ доказуемо
эквивалентна формуле из класса $\Sigma_{n+1}$, а также формуле
из класса~$\Pi_{n+1}$. В самом деле, если формула $\psi$ не имеет
параметра $\eta$, то она будет доказуемо эквивалентна формулам
$\exists\eta\,\psi$ и $\forall\eta\,\psi$ (одна импликация
является аксиомой, другая получается из $\psi\hm\to\psi$ по
правилу Бернайса). Таким образом, можно добавить фиктивный
квантор в начало кванторной приставки или в её конец; во втором
случае надо сослаться на лемму~2
раздела~\ref{renaming-variables} (с.~\pageref{changing-parts}).

        \item
Отрицание любой формулы из класса~$\Sigma_n$ доказуемо
эквивалентно некоторой формуле из класса~$\Pi_n$ и наоборот. В
самом деле, мы видели, что $\lnot\exists\xi\,\psi$ выводимо
эквивалентно $\forall\xi\,\lnot\psi$ и наоборот
(с.~\pageref{quantifier-duality}), так что отрицание можно
проносить внутрь, меняя по ходу дела кванторы на двойственные.

        \item
Конъюнкция любых двух формул из $\Pi_1$ доказуемо
эквивалентна некоторой формуле из~$\Pi_1$. Например,
конъюнкция $\forall x\,A(x)\hm\land \forall y\,B(y)$
доказуемо эквивалентна формуле ${\forall x \forall y\,
(A(x)\land B(y))}$. В самом деле, используя аксиомы про квантор
всеобщности, можно из $\forall x\, A(x)$ вывести $A(x)$, а из
$\forall y\, B(y)$ вывести $B(y)$, поэтому из их конъюнкции
выводится $A(x)\hm\land B(y)$, после чего можно навесить два
квантора всеобщности. В другую сторону:
выводим формулу ${\forall x\forall y\, (A(x)\land B(y))}\hm\to
A(x)$, затем применяем правило Бернайса \итд

\begin{problem}
        %
Покажите, что можно сэкономить один квантор и
использовать формулу ${\forall x\, (A(x)\land B(x))}$.
        %
\end{problem}

Общее рассуждение (для любых двух формул из класса $\Pi_1$)
почти столь же просто, надо лишь переименовать связанные
переменные, пользуясь теоремой~\ref{renaming-variables-theorem}
(с.~\pageref{renaming-variables-theorem}).
        \item
Аналогично можно доказать, что дизъюнкция двух формул
из класса~$\Sigma_1$ доказуемо эквивалентна некоторой формуле
класса~$\Sigma_1$. (Можно также перейти к двойственному
классу~$\Pi_1$, воспользовавшись уже известными свойствами
отрицания.)
        \item
Покажем теперь, что конъюнкция двух формул из класса $\Sigma_1$
доказуемо эквивалентна формуле класса~$\Sigma_1$ и что
дизъюнкция двух формул класса~$\Pi_1$ доказуемо эквивалентна
формуле класса~$\Pi_1$. Для этого надо воспользоваться
эквивалентностями вида
        $$
\exists x\,A(x) \land
\exists y\,B(y) \leftrightarrow
\exists x \exists y (A(x)\land B(y))
        $$
и
        $$
\forall x\,A(x) \lor
\forall y\,B(y) \leftrightarrow
\forall x \forall y (A(x)\lor B(y)).
        $$
Отметим кстати, что в них уже нельзя сэкономить квантор;
например, формула $\exists x (A(x)\land B(x))$ не равносильна
формуле $\exists x\, A(x)\hm\land \exists x\, B(x)$. (В одном из
выступлений времён начала перестройки
М.\,С.\,Гор\-ба\-чёв\glossary{\Горбачёв} сказал, что нужны \лк
преданные делу социализма, но квалифицированные специалисты\пк\т
впрочем, в газетной публикации \лк но\пк\ было заменено на
нейтральное \лк и\пк. Так вот, их существование не вытекает из
отдельного существования тех и других.)

Указанные эквивалентности, как легко видеть, общезначимы и
потому выводимы. Это совсем просто понять для первой из них
(чтобы найти пару объектов с заданными свойствами, надо найти
отдельно первый и второй члены пары). Вторая эквивалентность
немного сложнее\т проще всего заметить, что она переходит в
первую при добавлении отрицания. (Большая сложность отражает тот
факт, что вторая эквивалентность, в отличие от первой, не
является интуиционистски верной.)

        \item
        \label{sigma-pi-closed}%
Теперь легко понять, что конъюнкция и дизъюнкция двух формул
из класса $\Sigma_n$ (или $\Pi_n)$ доказуемо эквивалентны
формулам из того же класса. В самом деле, с помощью указанных
выше эквивалентностей можно слить кванторные приставки.
Например, формула
        $$
\forall x \exists y \, A(x,y) \lor \forall u \exists v \, B(u,v)
        $$
доказуемо эквивалентна сначала формуле
        $$
\forall x \forall u \, (\exists y\, A(x,y) \lor \exists v\, B(u,v)),
        $$
а затем формуле
        $$
\forall x \forall u \exists y \exists v\, (A(x,y)\lor B(u,v)).
        $$

\begin{problem}
        %
Как сэкономить один квантор в этом преобразовании?
        %
\end{problem}

\end{itemize}

Теперь всё готово для доказательства упомянутого в начале раздела
результата.

\begin{theorem}[о предварённой нормальной форме]
        \index{Теорема!о предварённой нормальной форме}%
Любая формула
произвольной сигнатуры доказуемо эквивалентна некоторой
формуле той же сигнатуры, имеющей предварённую нормальную
форму.
        %
\end{theorem}

\begin{proof}
        %
Индукция по построению формулы. Для атомарных формул это
очевидно. Отрицание переводит формулу класса~$\Sigma_n$ в
класс~$\Pi_n$ и наоборот. Конъюнкция и дизъюнкция: приведём
каждую формулу к предварённой нормальной форме, затем добавим
фиктивные кванторы так, чтобы они попали в один класс, а затем
воспользуемся доказанным утверждением. Импликация сводится к
дизъюнкции и отрицанию ($\varphi\hm\to\psi$ доказуемо
эквивалентно $\lnot\varphi\lor\psi$).
        %
\end{proof}

Отметим, что ни формулировка, ни доказательство этой теоремы не
предполагают замкнутости формулы.

\begin{problem}
        %
Приведите к предварённой нормальной форме формулу
$\forall x\,A(x)\hm\to\forall x\, B(x)$.
        %
\end{problem}

\begin{problem}
        %
Формулы~$\varphi$ и~$\psi$ принадлежат классу~$\Sigma_n$.
Найдём формулу в предварённой нормальной форме, выводимо
эквивалентную формуле $\varphi\hm\to\psi$. В каком классе
она окажется? (Указание: возможны разные варианты.)
        %
\end{problem}

\begin{problem}
        %
Применим описанный метод приведения к общезначимой формуле $\exists
x\forall y\, A(x,y)\hm\to\forall y\exists x A(x,y)$. Какая
предварённая формула получится? (Естественно, она будет
общезначимой.)
        %
\end{problem}

\section{Теорема Эрбрана}

Естественно ожидать, что вопрос о выводимости (или
общезначимости) формулы тем сложнее, чем сложнее сама формула. В
этом разделе (а также в следующем) мы рассмотрим его для формул
класса~$\Sigma_n$ и $\Pi_n$.

Начнём с самого простого случая. Пусть
$\varphi$\т бескванторная формула. Посмотрим, из каких атомарных
формул она составлена, и заменим их на пропозициональные
переменные (разные\т на разные, одинаковые\т на одинаковые).
Получится формула логики высказываний, которую мы будем называть
\emph{прототипом}%
\index{Прототип формулы|defin}%
\index{Формула!прототип|defin}
формулы~$\varphi$. Имеет место следующее (почти
очевидное) утверждение.

\begin{theorem}[выводимость бескванторных формул]
\index{Теорема!о выводимости бескванторных формул|defin}%
       \label{quantifier-free-provability}%
Бес\-кван\-тор\-ная формула выводима (общезначима) тогда и только
тогда, когда её прототип является тавтологией.
        %
\end{theorem}

\begin{proof}
        %
Если прототип формулы~$\varphi$ является тавтологией, то
формула~$\varphi$ является частным случаем пропозициональной
тавтологии и потому выводима и общезначима.

Пусть прототип формулы~$\varphi$ не является тавтологией. Можно
считать, что формула~$\varphi$ замкнута, поскольку свободные
переменные с точки зрения общезначимости и выводимости ничем не
отличаются от констант (мы уже отмечали это при доказательстве
теоремы о полноте). Построим интерпретацию, где формула
$\varphi$ будет ложной. Носителем её будут замкнутые термы.
Значения предикатов мы подберём так, чтобы атомарные формулы
принимали те самые значения, которые делают прототип
формулы~$\varphi$ ложным. Это возможно, так как значения
разных атомарных формул можно выбирать независимо (это значения
либо разных предикатов, либо одного и того же предиката, но на
разных термах).
        %
\end{proof}

Что можно сказать про общезначимость формул классов~$\Pi_1$
и~$\Sigma_1$? Для класса~$\Pi_1$ всё просто: общезначимость
формулы со свободными переменными равносильна общезначимости её
замыкания (которое получается, если навесить кванторы
всеобщности по всем переменным), поэтому формулы класса~$\Pi_1$
по существу ничем не отличаются от бескванторных.

Вопрос для класса~$\Sigma_1$ решается следующей теоремой:

\begin{theorem}[Эрбрана]
\index{Теорема!Эрбрана|defin}%
\glossary{\Эрбран}%
        \label{herbrand-theorem}%
Формула $\exists \xi_1\ldots\exists\xi_k\,\varphi$ (где
формула~$\varphi$\т бескванторная) общезначима тогда и только
тогда, когда найдётся конечный список подстановок
        \begin{gather*}
\varphi(t_1/\xi_1,\dots, t_k/\xi_k),\\
\varphi(u_1/\xi_1,\dots, u_k/\xi_k),\\
\myboxfil{\varphi(u_1/\xi_1,\dots, u_k/\xi_k),}\,\\
\varphi(w_1/\xi_1,\dots, w_k/\xi_k)
        \end{gather*}
(вместо переменных подставляются термы нашей сигнатуры),
дизъюнкция которых общезначима.
        %
\end{theorem}

Заметим, что дизъюнкция, о которой идёт речь в теореме, является
бескванторной формулой. По
теореме~\ref{quantifier-free-provability} она общезначима тогда
и только тогда, когда является частным случае пропозициональной
тавтологии.

Прежде чем доказывать эту теорему, приведём пример. Рассмотрим
формулу
        $$
\exists x \,( A(c,x)\to A(x,d) )
        $$
(в которой $c$ и~$d$\т константы). Она общезначима; соответствующий
набор состоит из подстановок $c/x$ и $d/x$. В самом деле, формула
        $$
(A(c,c)\to A(c,d))\lor (A(c,d)\to A(d,d))
        $$
истинна как при истинном~$A(c,d)$, так и при ложном.
Заметим, что в этом примере нам понадобились две подстановки.

\begin{proof}
        %
Доказательство теоремы Эрбрана. В одну сторону утверждение
очевидно: если общезначима дизъюнкция подстановок, то
общезначима формула с квантором. (Мы уже использовали это
при элиминации кванторов в
разделе~\ref{quantifier-elimination} при доказательстве
теоремы~\ref{quantifier-elimination-integers-successor}.)

Докажем обратное утверждение. Будем считать, что
формула~$\varphi$ не содержит переменных,
кроме~$\xi_1,\dots,\xi_k$ (как мы уже замечали, остальные
переменные можно заменить константами). Рассмотрим (бесконечное)
множество формул
        $$
\lnot\varphi(t_1/\xi_1,\dots,t_k/\xi_k)
        $$
для всевозможных наборов замкнутых термов $t_1,\dots,t_k$. Если
это множество противоречиво, всё доказано (тогда выводима
дизъюнкция подстановок, отрицания которых
используются при выводе противоречия).
Если оно
непротиворечиво, то существует интерпретация, в которой все эти
формулы истинны. Мы не можем утверждать, что в этой интерпретации
ложна формула
        $$
\exists\xi_1\dots\exists\xi_k\,\varphi
        $$
(носитель интерпретации может содержать элементы, не являющиеся
значениями замкнутых термов). Однако если мы выбросим лишние
элементы и оставим только значения термов, то эта формула
станет ложной, так что она не общезначима.
        %
\end{proof}

Теорему Эрбрана можно сформулировать чисто
синтаксически: если выводима $\Sigma_1$\д формула $\exists
\xi_1\ldots\exists \xi_k \varphi$, то можно найти конечное число
подстановок, дизъюнкция которых выводима. Можно предложить и
доказательство, не использующее понятия
общезначимости. Такое
доказательство приведено, например, в книге
Клини%
\glossary{\Клини}%
~\cite{kleene-logic} (для генценовского%
        \glossary{\Генцен}
варианта исчисления предикатов) и в книге Шёнфилда%
        \glossary{\Шёнфилд}%
~\cite{shoenfield} (для
        \glossary{\Гильберт}%
гильбертовского варианта). Синтаксическое доказательство (в
отличие от нашего) конструктивно: по выводу $\Sigma_1$\д формулы
можно алгоритмически указать соответствующие термы.

Если сигнатура не содержит функциональных символов, то с помощью теоремы
Эрбрана можно алгоритмически проверять выводимость формул
класса $\Sigma_1$, поскольку число возможных подстановок
конечно. Это же можно сказать и про формулы класса~$\Pi_2$, так
как внешние кванторы всеобщности можно отбросить, не меняя
выводимости.

Естественный вопрос: можно ли построить аналогичные алгоритмы
для следующих классов? Отрицательный ответ даётся в следующем
разделе.

\section{Сколемовские функции}
\index{Сколемовская функция|defin}%
\index{Функция!сколемовская|defin}%

В этом разделе мы в разных формах используем ровно одну идею:
истинность утверждения
        $$
\forall x \exists y \, A(x,y)
        $$
равносильна существованию функции, которая по любому~$x$ даёт такой~$y$, что
$A(x,y)$. Это утверждение нельзя записать в виде эквивалентности
        $$
\forall x\exists y \, A(x,y) \Leftrightarrow \exists f \forall x\, A(x,f(x)),
        $$
поскольку в нашем языке нет квантора по функциям
и~$\exists f$ мы писать не имеем права. (Языки, содержащие кванторы по
множествам и функциям, называются языками \emph{второго
порядка}%
\index{Языки!второго порядка|defin}
и мы их не рассматриваем.)

Тем не менее это утверждение можно сформулировать и в наших
терминах. Пусть, например, имеется формула $\varphi$ с двумя
параметрами $x$ и $y$. Тогда замкнутая формула
        $\forall x \exists y\,\varphi$
выполнима тогда и только тогда, когда выполнима формула
        $\forall x \varphi(f(x)/y)$,
где $f$\т новый одноместный функциональный символ. (Аккуратный
читатель поправит: надо ещё требовать, чтобы подстановка~$f(x)$
вместо~$y$ была корректна. Мы уже знаем, что переменные можно
переименовывать, поэтому будем легкомысленно считать, что все
необходимые переименования уже сделаны.)

Аналогичное можно преобразовать произвольные предварённые
формулы. Например, формула
        $$
\forall x \forall y \exists z \forall u \exists v\,\varphi (x,y,z,u,v)
        $$
выполнима тогда и только тогда, когда выполнима формула
        $$
\forall x \forall y \forall u\,\varphi (x,y,f(x,y),u,g(x,y,u))
        $$
\label{skolem-functions-example}%
(здесь $\varphi(x,y,z,u,v)$\т формула, не имеющая
параметров, кроме явно указанных $x,y,z,u,v$, запись
$\varphi(x,y,f(x,y),u,g(x,y,u))$
обозначает результат соответствующих подстановок, которые мы
предполагаем корректными, а $f$ и~$g$\т функциональные символы,
не встречающиеся в формуле~$\varphi$).

Сходное преобразование имеют в виду преподаватели
математического анализа, которые иногда записывают определение
предела ($\forall \varepsilon \exists\delta\ldots$) в несколько
странной для логика форме $\forall \varepsilon \exists
\delta=\delta(\varepsilon)\ldots$\т имеется в виду, что если для
каждого~$\varepsilon$ найдётся~$\delta$, то это самое~$\delta$
представляет собой функцию от~$\varepsilon$.

Отметим, что это рассуждение использует аксиому выбора, когда из
различных возможных (для данного~$\varepsilon$) значений
$\delta$ мы выбираем какое-то одно и объявляем его значением
функции $\delta(\varepsilon)$.

\begin{problem}
        %
Казалось бы, выбор~$v$ в приведённом выше примере зависит от
$x,y,z,u$, так что следовало бы написать
$\varphi(x,y,f(x,y),u,g(x,y,f(x,y),u))$, но мы так не делаем.
Почему это допустимо?
        %
\end{problem}

В общем случае мы получаем такое утверждение:

\begin{theorem}
        \label{skolem-functions-1}
Для всякой замкнутой формулы~$\tau$ сигнатуры~$\sigma$ можно
указать формулу~$\tau'$ класса $\Pi_1$ сигнатуры~$\sigma$ с
добавленными функциональными символами, которая выполнима или
невыполнима одновременно с формулой~$\tau$. При этом
преобразование $\tau\mapsto\tau'$ эффективно (выполняется
некоторым алгоритмом).
        %
\end{theorem}

\begin{proof}
        %
Приведя $\tau$ к предварённой нормальной форме, получим
доказуемо эквивалентную формулу (выполнимую или невыполнимую
одновременно с~$\tau$). После этого применяем
описанное выше преобразование.
        %
\end{proof}

Формула $\tau$ невыполнима тогда и только тогда, когда её
отрицание общезначимо. Поэтому наши рассуждения показывают, что,
скажем, формула $\lnot \forall x \exists y\,\psi(x,y)$
общезначима одновременно с формулой $\lnot \forall x
\,\psi(x,f(x))$. Внося отрицание внутрь и заменяя $\lnot\psi$ на
$\varphi$, получаем такое утверждение: формулы
        $$
\exists x \forall y\,\varphi(x,y)
        \quad \text{и}\quad
\exists x \,\varphi(x,f(x))
        $$
одновременно общезначимы. (Это утверждение чуть менее наглядно,
чем двойственное ему утверждение о выполнимости.) В общем
виде двойственное к теореме~\ref{skolem-functions-1}
утверждение выглядит так:

\begin{theorem}
        \label{skolem-functions-2}
Для всякой замкнутой формулы~$\tau$ сигнатуры~$\sigma$ можно
указать формулу~$\tau'$ класса $\Sigma_1$ сигнатуры~$\sigma$ с
добавленными функциональными символами, которая общезначима или
необщезначима одновременно с формулой~$\tau$. При этом
преобразование $\tau\mapsto\tau'$ эффективно (выполняется
некоторым алгоритмом).
        %
\end{theorem}

Заметим, что к формуле $\tau'$ можно применить теорему Эрбрана
(с.~\pageref{herbrand-theorem}): она общезначима тогда и только тогда,
когда дизъюнкция нескольких подстановок является тавтологией.

Теорема о полноте позволяет заменить в этой формулировке
общезначимость на выводимость: формулы $\tau$ и $\tau'$
одновременно выводимы. После этого естественно искать явный
метод, который преобразует вывод формулы~$\tau$ в вывод
формулы~$\tau'$ и обратно. Такой метод действительно существует,
но для этого требуется более детальный анализ структуры
выводов, при котором удобно пользоваться исчислениями
генценовского типа.

Идея использования функций вместо групп кванторов
$\forall\exists$ восходит к Эрбрану%
\glossary{\Эрбран}
и Сколему.%
\glossary{\Сколем}
Такие функции
иногда называют \лк эрбрановскими%
\index{Эрбрановская функция|defin}%
\index{Функция!эрбрановская|defin}%
\пк\ или \лк сколемовскими\пк,%
\index{Сколемовская функция|defin}%
\index{Функция!сколемовская|defin}
а их добавление \т \лк сколемизацией\пк.%
\index{Сколемизация|defin}
Есть также термин \лк
сколемовская нормальная форма\пк,%
\index{Сколемовская нормальная форма|defin}%
\index{Нормальная форма!сколемовская|defin}
но в ней добавляются не
функциональные символы, а предикатные, и получается формула не
класса $\Sigma_1$, как в теореме~\ref{skolem-functions-2}, а
класса~$\Sigma_2$.

\begin{theorem}[о сколемовской нормальной форме]
\index{Теорема!о сколемовской нормальной форме|defin}%
        \label{skolem-normal-form}%
Для всякой замкнутой формулы~$\tau$ сигнатуры~$\sigma$ можно
указать формулу~$\tau'$ класса $\Sigma_2$ сигнатуры~$\sigma$ с
добавленными предикатными символами, которая общезначима или
необщезначима одновременно с формулой~$\tau$. При этом
преобразование $\tau\mapsto\tau'$ эффективно (выполняется
некоторым алгоритмом).
        %
\end{theorem}

\begin{proof}
        %
Как и раньше, нам будет удобнее говорить о выполнимости и
доказывать, что для всякой формулы~$\tau$ найдётся одновременно
с ней выполнимая формула~$\tau'$ из класса $\Pi_2$.
Построение такой формулы мы объясним на примере. Пусть
исходная формула имеет вид
        $$
\forall x \forall y \exists z \forall u \exists v\,\varphi(x,y,z,u,v).
        $$
Мы теперь не можем ввести функции~$f$ и~$g$, как это делалось
выше, на с.~\pageref{skolem-functions-example}. Поэтому мы введём
предикаты~$F$ и~$G$, заменяющие графики этих функций, и напишем
формулу
       \begin{align*}
\forall x\forall y \exists z&\,\, F(x,y,z)
        \land \\
\forall x \forall y \forall u \exists v&\,\, G(x,y,u,v)
        \land\\
\forall x \forall y \forall z \forall u \forall v&\,
        (F(x,y,z)\land G(x,y,u,v) \to \varphi(x,y,z,u,v))
        \end{align*}
Если исходная формула выполнима, то новая тоже выполнима:
достаточно взять в качестве $F$ и $G$ графики сколемовских
функций. Напротив, если новая формула выполнима, то выполнима и
старая (более того, из новой формулы следует старая): надо взять
$z$ и $v$ согласно первым двум строкам и заметить, что согласно
третьей строке они подойдут.

Такая конструкция применима к любой предварённой форме и даёт
конъюнкцию $\Pi_2$\д формул (последняя из которых будет
даже и $\Pi_1$\д формулой). А мы знаем~(с.~\pageref{sigma-pi-closed}),
что такая конъюнкция эквивалентна
$\Pi_2$\д формуле.
        %
\end{proof}

\begin{problem}
        %
Дайте синтаксическое доказательство теоремы о сколемовской
нормальной форме (показав, что из выводимости формулы следует
выводимость её сколемовской нормальной формы и наоборот).
(Указание: это проще, чем для формул с функциональными
символами, и не требуется использовать генценовское
исчисление.)
        %
\end{problem}

Утверждения этого раздела сводят вопрос о выводимости
произвольной формулы исчисления предикатов к выводимости
$\Sigma_1$\д формулы (с функциональными символами). Если мы
запрещаем функциональные символы, то вопрос о выводимости
произвольной формулы сводится к выводимости $\Sigma_2$\д формулы.

Как мы увидим в разделе~\ref{noncomplete-undecidable} (теорема
Чёрча%
 \index{Теорема!Чёрча}%
 \glossary{\Чёрч}, c.~\pageref{church-theorem}),
вопрос о выводимости произвольных формул языка первого порядка
неразрешим: не существует алгоритма, который бы по произвольной
замкнутой формуле определял бы, выводима она или нет. Результаты
этого раздела показывают, что уже для формул класса~$\Sigma_1$
(с функциональными символами) или~$\Sigma_2$ (без них) такого
алгоритма не существует, поскольку из него можно было бы
получить и общий алгоритм. (В~предыдущем разделе мы видели, что
для формул класса~$\Sigma_1$ без функциональных символов такой
алгоритм существует.)
